\documentclass[10pt, twocolumn]{article}
\usepackage[utf8]{inputenc}
\usepackage{kotex}
\usepackage{amsmath, amssymb, amsfonts}
\usepackage{graphicx}
\usepackage{booktabs}
\usepackage{multirow}
\usepackage{algorithm}
\usepackage{algpseudocode}
\usepackage[hidelinks]{hyperref}
\usepackage{cite}
\usepackage{geometry}
\geometry{a4paper, margin=2cm}

\title{\textbf{강화학습 기반 5G 네트워크 슬라이싱의 \\Three-Part Tariff 동적 가격 최적화}}

\author{
  \textit{저자명} \\
  소속기관 \\
  \texttt{email@example.com}
}

\date{}

\begin{document}
\maketitle

% ============================================================
% 초록
% ============================================================
\begin{abstract}
5G 네트워크 슬라이싱은 URLLC(초저지연 고신뢰)와 eMBB(고대역폭) 등 이질적 서비스를 단일 물리 인프라 위에서 제공한다.
본 연구에서는 서비스 제공자의 수익을 극대화하기 위해 Three-Part Tariff(3PT) 구조---고정요금~$F$, 허용량~$\bar{Q}$, 초과요금~$p$---의 동적 가격 결정 문제를 마르코프 결정 과정(MDP)으로 정식화하고, Soft Actor-Critic(SAC) 강화학습으로 해결하는 프레임워크를 제안한다.
환경 모델은 로그정규 트래픽 분포\cite{alasmar2021}, 로지스틱 이탈 모델\cite{ahn2006, ribeiro2023}, 가격 탄력적 Poisson 유입\cite{chen2022}, 그리고 3GPP TS~22.261 기반 QoS 패널티\cite{kim2019}로 구성되며, 모든 구성요소가 기존 문헌에 근거한다.
시뮬레이션을 통해 SAC 기반 동적 가격 정책이 고정 가격 및 단순 요금제 대비 우수한 순수익을 달성함을 검증한다.
\end{abstract}

\textbf{키워드}: 5G 네트워크 슬라이싱, Three-Part Tariff, 강화학습, 동적 가격 최적화, Soft Actor-Critic

% ============================================================
% 1. 서론
% ============================================================
\section{서론}

5G 네트워크 슬라이싱은 단일 물리적 네트워크 인프라를 다수의 논리적 네트워크(슬라이스)로 분할하여, 각 슬라이스가 차별화된 서비스 품질(QoS)을 보장하는 핵심 기술이다\cite{kim2019}.
URLLC(Ultra-Reliable Low-Latency Communication) 슬라이스는 자율주행, 원격수술 등에 1\,ms 이내 지연과 99.999\% 신뢰성을 요구하며, eMBB(enhanced Mobile Broadband) 슬라이스는 영상 스트리밍, 클라우드 게임 등에 100\,Mbps 이상의 처리량을 목표로 한다\cite{3gpp22261}.
5G 네트워크 슬라이싱 시장은 2025년 8.4억 달러에서 2030년 50.7억 달러로 연평균 43.3\% 성장이 전망된다.

이러한 이질적 슬라이스에 대한 \textbf{가격 설계}는 서비스 제공자의 수익과 가입자 유지에 직결되는 핵심 과제이다.
실무에서 통신 요금제는 주로 \textbf{Three-Part Tariff(3PT)}---고정 접속료~$F$, 무료 허용량~$\bar{Q}$, 초과 단위당 요금~$p$---구조를 따르며\cite{wu2023}, 이는 Two-Part Tariff(2PT)이나 정액제보다 비용 효율적임이 이론적으로 입증되었다\cite{wu2023}.
그러나 기존 3PT 최적화 연구\cite{wu2023}는 정적 수요를 가정하여, 가입자 이탈·유입의 동적 변화와 확률적 QoS 변동을 고려하지 못한다.

본 연구에서는 이 한계를 극복하기 위해 다음 세 가지 기여를 제시한다:
\begin{itemize}
  \item \textbf{C1}: URLLC/eMBB 다중 슬라이스에 대한 3PT 동적 가격 결정 문제의 MDP 정식화
  \item \textbf{C2}: 로그정규 트래픽\cite{alasmar2021}, 로지스틱 이탈\cite{ahn2006}, Poisson 유입\cite{chen2022}, QoS 패널티\cite{kim2019} 등 모든 구성요소를 기존 문헌에 근거한 환경 모델
  \item \textbf{C3}: SAC 및 PPO 강화학습과 다수 베이스라인 간 실증 비교
\end{itemize}

% ============================================================
% 2. 시스템 모델
% ============================================================
\section{시스템 모델}

\subsection{네트워크 슬라이싱 구조}

서비스 제공자는 슬라이스 유형 $s \in \mathcal{S} = \{U, E\}$ (URLLC, eMBB)를 운영한다.
시간은 이산 스텝 $t \in \{1, 2, \ldots, T\}$로 구성되며, 한 스텝은 1시간에 해당한다 ($T=720$, 1개월).
각 슬라이스~$s$의 시각~$t$ 시작 시 가입자 수를 $N_{s,t} \in \mathbb{Z}^+$로 표기한다.

\subsection{Three-Part Tariff 청구 구조}

사용자 $i$의 슬라이스~$s$에서 시각~$t$의 청구액은\cite{wu2023}:
\begin{equation}
  B_{i,s,t} = F_{s,t} + \max\!\big(0,\; q_{i,s,t} - \bar{Q}_s\big) \cdot p_{s,t}
  \label{eq:bill}
\end{equation}
여기서 $F_{s,t} \in \mathbb{R}^+$는 고정요금(학습), $\bar{Q}_s \in \mathbb{R}^+$는 허용량(고정 상수), $p_{s,t} \in \mathbb{R}^+$는 초과요금(학습), $q_{i,s,t}$는 데이터 사용량이다.

\subsection{데이터 사용량 모델}

개별 사용자의 시간당 데이터 사용량은 로그정규분포를 따른다\cite{alasmar2021}:
\begin{equation}
  q_{i,s,t} \sim \text{LogNormal}(\mu_s, \sigma_s^2)
  \label{eq:usage}
\end{equation}
즉 $\ln q_{i,s,t} \sim \mathcal{N}(\mu_s, \sigma_s^2)$이며, 기대값과 분산은 다음과 같다:
\[
  \mathbb{E}[q] = \exp\!\Big(\mu_s + \frac{\sigma_s^2}{2}\Big), \quad
  \text{Var}[q] = \big[\exp(\sigma_s^2)-1\big]\exp(2\mu_s + \sigma_s^2)
\]

이에 따른 기대 초과 사용량의 폐쇄형 해는:
\begin{equation}
  \mathbb{E}\big[\max(0,\, q-\bar{Q})\big] = \mathbb{E}[q]\,\Phi(d_1) - \bar{Q}\,\Phi(d_2)
  \label{eq:overage}
\end{equation}
여기서 $d_1 = (\mu_s + \sigma_s^2 - \ln\bar{Q}_s)/\sigma_s$, $d_2 = d_1 - \sigma_s$, $\Phi$는 표준정규 CDF이다.

\subsection{가입자 이탈 모델}

기존 가입자의 이탈 확률은 로지스틱 함수로 모델링한다\cite{ahn2006, ribeiro2023}:
\begin{equation}
  P_{s,t}^{\text{dep}} = \sigma\!\Big(\gamma_{0,s} + \gamma_F \tilde{F}_{s,t} + \gamma_p \tilde{p}_{s,t} - \gamma_{\eta,s}\, \eta_{s,t-1}\Big)
  \label{eq:depart}
\end{equation}
여기서 $\sigma(x) = 1/(1+e^{-x})$, $\tilde{F}_{s,t} = F_{s,t}/F_{\text{ref},s}$, $\tilde{p}_{s,t} = p_{s,t}/p_{\text{ref},s}$는 정규화 가격, $\eta_{s,t-1}$은 이전 시간의 QoS이다.
$\gamma_{0,s} < 0$은 기본 유지 경향, $\gamma_F, \gamma_p > 0$은 가격 민감도, $\gamma_{\eta,s} > 0$은 QoS 민감도를 나타낸다.
특히 URLLC 사용자의 QoS 이탈 민감도가 eMBB의 6배이다 ($\gamma_{\eta,U} = 3.0 \gg \gamma_{\eta,E} = 0.5$).
각 가입자는 독립적으로 이탈한다: $\text{Leave}_{i,s,t} \sim \text{Bernoulli}(P_{s,t}^{\text{dep}})$.

\subsection{가입자 유입 모델}

신규 가입 확률은 현재 가격에만 의존한다\cite{chen2022} (QoS 경험 없음):
\begin{equation}
  P_{s,t}^{\text{arr}} = \sigma\!\Big(\beta_{0,s} - \beta_F \tilde{F}_{s,t} - \beta_p \tilde{p}_{s,t}\Big)
  \label{eq:arrive}
\end{equation}
신규 가입자 수는 조절된 Poisson 과정을 따른다:
\begin{equation}
  N_{s,t}^{\text{new}} \sim \text{Poisson}\!\big(\lambda_{\max,s} \cdot P_{s,t}^{\text{arr}}\big)
  \label{eq:poisson}
\end{equation}
활성 사용자 수는 $N_{s,t}^{\text{active}} = (N_{s,t} - N_{s,t}^{\text{leave}}) + N_{s,t}^{\text{new}}$이다.

\subsection{QoS 및 보상}

QoS 만족도는 외생적 확률 변수로 모델링한다:
\begin{equation}
  \eta_{s,t} \sim \text{Uniform}(\eta_s^{\text{low}},\, \eta_s^{\text{high}})
  \label{eq:qos}
\end{equation}

보상은 총 수익에서 SLA 위반 패널티를 감산한다\cite{kim2019}:
\begin{align}
  R_t &= \sum_{s \in \{U,E\}} \sum_{i \in \mathcal{A}_{s,t}} B_{i,s,t} \label{eq:revenue} \\
  \text{Pen}_t &= \sum_{s \in \{U,E\}} w_s \cdot N_{s,t}^{\text{active}} \cdot \max\!\big(0,\, \eta_s^{\text{tgt}} - \eta_{s,t}\big) \label{eq:penalty} \\
  r_t &= R_t - \text{Pen}_t \label{eq:reward}
\end{align}
여기서 $w_s$는 SLA 패널티 가중치이며, URLLC의 엄격한 QoS 요구를 반영하여 $w_U \gg w_E$로 설정한다.

% ============================================================
% 3. MDP 정식화 및 강화학습
% ============================================================
\section{MDP 정식화 및 강화학습}

\subsection{MDP 정의}

상기 시스템을 MDP 튜플 $(\mathcal{S}, \mathcal{A}, P, r, \gamma)$로 정식화한다:

\begin{table}[h]
\centering
\caption{MDP 구성요소}
\label{tab:mdp}
\begin{tabular}{@{}ll@{}}
\toprule
구성요소 & 정의 \\
\midrule
상태 & $\mathcal{S}_t = (N_{U,t},\, N_{E,t},\, \eta_{U,t-1},\, \eta_{E,t-1})$ --- 4차원 \\
액션 & $a_t = (F_{U,t},\, p_{U,t},\, F_{E,t},\, p_{E,t}) \in \mathcal{A}$ --- 4차원 연속 \\
보상 & $r_t = R_t - \text{Pen}_t$ \\
전이 & $N_{s,t+1} = N_{s,t}^{\text{active}}$, QoS 이월 \\
목표 & $\max_\pi \mathbb{E}\!\Big[\sum_{t=1}^{T} \gamma^{t-1} r_t\Big],\;\; \gamma \in [0,1)$ \\
\bottomrule
\end{tabular}
\end{table}

상태에 이전 시간 QoS~$\eta_{s,t-1}$을 포함하는 이유는, 이탈 확률(\ref{eq:depart})이 이전 QoS에 의존하기 때문이다\cite{ahn2006}.
에이전트가 QoS를 직접 제어할 수는 없으나, QoS 변동을 관찰하여 \textbf{반응적 가격 조정}(예: QoS 저하 시 이탈 방지를 위한 가격 인하)이 가능하다.

\subsection{Soft Actor-Critic (SAC)}

본 연구의 주력 알고리즘으로 SAC\cite{haarnoja2018}를 채택한다.
SAC는 off-policy 방식으로, 엔트로피 정규화된 목적함수를 최대화한다:
\begin{equation}
  \max_\phi \;\mathbb{E}\!\Big[Q_\theta(s, a) - \alpha \log \pi_\phi(a|s)\Big]
  \label{eq:sac}
\end{equation}
여기서 $\alpha$는 자동 조정되는 온도 파라미터이다.
Twin Q-네트워크($Q_{\theta_1}, Q_{\theta_2}$)로 과대평가를 방지하며, 연속 4차원 액션 공간에서의 탐색-활용 균형을 자동으로 조율한다.

Actor와 Critic 모두 3층 MLP (256-256-256, ReLU)로 구성하며, 출력층은 시그모이드 활성화 후 액션 범위 $F \in [0, 100]$, $p \in [0, 20]$으로 스케일링한다.

\subsection{비교 알고리즘: PPO}

On-policy 비교를 위해 PPO (Proximal Policy Optimization)를 사용한다.
Clipped surrogate 목적함수로 정책 변화를 제한하여 안정적 학습을 보장한다.

\subsection{베이스라인 정책}

학습된 RL 정책의 우수성을 검증하기 위해 6개 베이스라인과 비교한다:
(1)~Static-Oracle (그리드 서치 최적 고정 가격),
(2)~Static-Heuristic (기준가격 고정),
(3)~2PT-SAC ($\bar{Q}=0$으로 설정한 SAC),
(4)~Flat-SAC ($p=0$으로 고정, $F$만 학습),
(5)~Myopic-SAC ($\gamma=0$),
(6)~Random (균등 랜덤 액션).

% ============================================================
% 참고문헌
% ============================================================
\begin{thebibliography}{9}

\bibitem{wu2023}
J.~F. Wu and B.~Behzad,
``Optimal three-part tariff pricing with Spence-Mirrlees reservation prices,''
\textit{Math. Methods Oper. Res.}, vol.~97, pp.~233--258, 2023.

\bibitem{alasmar2021}
M.~Alasmar, R.~Clegg, N.~Zakhleniuk, and G.~Pavlou,
``Internet traffic volumes are not Gaussian---they are log-normal: An 18-year longitudinal study,''
\textit{IEEE/ACM Trans. Netw.}, vol.~29, no.~3, pp.~1266--1279, 2021.

\bibitem{ahn2006}
J.-H. Ahn, S.-P. Han, and Y.-S. Lee,
``Customer churn analysis: Churn determinants and mediation effects of partial defection in the Korean mobile telecommunications service industry,''
\textit{Telecom. Policy}, vol.~30, pp.~552--568, 2006.

\bibitem{chen2022}
J.~Chen, L.~Jiang, and S.~A.~S. Shah,
``An empirical model of the effects of `bill shock' regulation in mobile telecommunication markets,''
\textit{Int. J. Ind. Organ.}, vol.~83, p.~102848, 2022.

\bibitem{ribeiro2023}
H.~Ribeiro, B.~Barbosa, A.~C. Moreira, and R.~G. Rodrigues,
``Determinants of churn in telecommunication services: A systematic literature review,''
\textit{Mgmt Rev. Quarterly}, vol.~74, pp.~1327--1364, 2023.

\bibitem{3gpp22261}
3GPP,
``Service requirements for the 5G system,''
3GPP TS 22.261, 2023.

\bibitem{kim2019}
Y.~Kim, S.~Kim, and H.~Lim,
``Reinforcement learning based resource management for network slicing,''
\textit{Appl. Sci.}, vol.~9, no.~11, p.~2361, 2019.

\bibitem{haarnoja2018}
T.~Haarnoja, A.~Zhou, P.~Abbeel, and S.~Levine,
``Soft actor-critic: Off-policy maximum entropy deep reinforcement learning with a stochastic actor,''
in \textit{Proc. ICML}, 2018.

\end{thebibliography}

\end{document}
